\documentclass[a4paper,12pt]{article}
\usepackage{cmap}
\usepackage[T2A]{fontenc}
\usepackage[utf8]{inputenc}
\usepackage[russian]{babel}
\usepackage[left=25mm,right=20mm,top=20mm,bottom=20mm]{geometry}
\usepackage{setspace,enumitem,longtable,array,booktabs,tikz,amsmath}
\usetikzlibrary{arrows.meta,positioning,shapes.geometric,fit,backgrounds}
\usepackage[unicode,hidelinks]{hyperref}
\onehalfspacing
\setlength{\parindent}{1.25cm}
\setlength{\parskip}{2pt}
\setlist[enumerate,itemize]{nosep,leftmargin=*,topsep=4pt}
\setlist[description]{nosep,topsep=4pt}
\newcolumntype{P}[1]{>{\raggedright\arraybackslash}p{#1}}
\setcounter{tocdepth}{1}
\setlength{\emergencystretch}{3em}
\newcommand{\field}[2]{\par\noindent\textbf{#1:} #2\par}
\newcommand{\actor}[3]{\draw (#1,#2+0.27) circle (0.15);\draw (#1,#2+0.12)--(#1,#2-0.35);\draw (#1-0.25,#2-0.05)--(#1+0.25,#2-0.05);\draw (#1,#2-0.35)--(#1-0.23,#2-0.65);\draw (#1,#2-0.35)--(#1+0.23,#2-0.65);\node[align=center,font=\small] at (#1,#2-1.0) {#3};}
\begin{document}
\pagenumbering{gobble}
\begin{titlepage}
\thispagestyle{empty}
\begin{center}
Федеральное государственное образовательное учреждение высшего образования

«Национальный исследовательский университет ИТМО»

\vspace{3.8cm}
{\Large Курсовая работа по дисциплине}

{\Large «Информационные системы»}

\vspace{0.6cm}
{\Large Этап 1. Анализ предметной области\par и проектирование информационной системы}

\vspace{0.6cm}
{\large Сервис бронирования предметов\par для жителей общежития на Вяземском переулке}

\vfill
\begin{flushright}
Выполнили:

Западовников Алексей Викторович. P3313

Кайгородова Александра Андреевна

Проверил:

Кобик Никита Алексеевич
\end{flushright}
\vfill
Санкт-Петербург 2026
\end{center}
\end{titlepage}
\pagestyle{empty}
\tableofcontents
\clearpage
\pagenumbering{arabic}
\setcounter{page}{3}
\pagestyle{plain}

\section{Проблема}
Жителям общежития на Вяземском переулке периодически нужны вещи для кратковременного использования: фен, утюг, инструменты или настольные игры. Покупать каждому отдельный экземпляр затратно, а хранить редко используемые вещи в жилых комнатах неудобно.

Общий фонд в отдельной комнате позволяет пользоваться такими предметами совместно. Однако без единого учёта трудно узнать, какие вещи доступны, кто ими пользуется и когда они вернутся. Договорённости в переписке могут теряться, а несколько жителей могут рассчитывать на один предмет одновременно. Нужна информационная система для ведения каталога, бронирования и учёта возврата.

\section{Решение и взаимодействие с системой}
Разрабатывается веб-сервис бронирования предметов общей комнаты. Принесённые вещи становятся общим имуществом; принёсший предмет житель пользуется им наравне с остальными. Каждая карточка соответствует отдельному экземпляру, который получает идентификатор и метку на самой вещи.

В системе четыре актора: \textbf{гость} просматривает каталог и регистрируется; \textbf{потенциальный пользователь} ожидает одобрения администратора и может только просматривать каталог и карточки, без повторной регистрации; \textbf{зарегистрированный пользователь} имеет подтверждённый доступ и пользуется фондом; \textbf{администратор} управляет доступом, каталогом и сообщениями о проблемах. Администратору также доступны функции зарегистрированного пользователя. Использование бесплатно и основано на доверии: получение и возврат пользователь отмечает самостоятельно, без подтверждения другим человеком.

\par\noindent\textbf{Основные сущности предметной области:}
\begin{itemize}
\item \textbf{Учётная запись} --- сведения о жителе, роль и состояние доступа: «Ожидает подтверждения», «Активен» или «Заблокирован». Потенциальный пользователь --- состояние этой записи; одобрение меняет его на «Активен», не создавая новую запись.
\item \textbf{Предмет} --- отдельная вещь общего фонда с идентификатором, описанием, состоянием и правилами использования.
\item \textbf{Категория} --- группа предметов одного назначения для поиска и фильтрации каталога.
\item \textbf{Бронь} --- резервирование одного предмета одним пользователем на заданный интервал; содержит статус и время получения.
\item \textbf{Возврат} --- запись о завершении использования по брони с временем возврата и обязательной фотографией предмета в общей комнате.
\item \textbf{Сообщение о проблеме} --- обращение пользователя по поводу предмета с описанием неисправности или отсутствия и решением администратора.
\end{itemize}

\clearpage
Порядок взаимодействия:
\begin{enumerate}
\item Гость создаёт учётную запись и становится потенциальным пользователем. До решения администратора ему доступен только каталог; повторная регистрация недоступна. После подтверждения он становится зарегистрированным пользователем и получает доступ к бронированию и пополнению фонда.
\item Для пополнения фонда приносит вещь в общую комнату, создаёт карточку с описанием, состоянием и правилами использования, маркирует предмет.
\item Для использования находит предмет в каталоге и выбирает интервал. Система автоматически бронирует свободное время.
\item В свой интервал берёт вещь из общей комнаты и нажимает «Я взял вещь».
\item Возвращает предмет в ту же комнату, нажимает «Я вернул вещь» и прикрепляет фотографию предмета на месте хранения. Система завершает использование.
\end{enumerate}

\clearpage
\section{Функциональные и нефункциональные требования}
\begingroup
\small
\setstretch{1.03}
\subsection{F. Функциональные требования}
Акторы: гость, потенциальный пользователь, зарегистрированный пользователь и администратор.

\par\noindent\textbf{Гость должен иметь возможность:}
\begin{itemize}[itemsep=0pt]
\item \textbf{FR-01.} Создать учётную запись, указав имя, уникальный логин и пароль; начальное состояние доступа --- «Ожидает подтверждения». Повторная регистрация для существующей записи запрещена; занятый логин отклоняется независимо от состояния доступа.
\item \textbf{FR-02.} Просмотреть каталог и карточки предметов без регистрации, без доступа к истории пользователей и фотографиям возврата.
\end{itemize}
\par\noindent\textbf{Потенциальный пользователь:} может только просматривать каталог и карточки по FR-02, без истории пользователей и фотографий возврата. Создание новой учётной записи, бронирование, пополнение фонда и остальные операции пользователя ему запрещены. Вход в систему до подтверждения запрещён; каталог доступен без авторизации. Регистрация не открывает сессию учётной записи.

\par\noindent\textbf{Зарегистрированный пользователь с подтверждённым доступом должен иметь возможность:}
\begin{itemize}[itemsep=0pt]
\item \textbf{FR-03.} Войти в систему по логину и паролю. Вход доступен только после подтверждения администратором. Для записи «Ожидает подтверждения» система отклоняет вход без создания сессии. После блокировки ранее подтверждённой записи вход сохраняется только с ограничениями FR-28.
\item \textbf{FR-04.} Выйти из системы с завершением своей сессии; это доступно для любой открытой сессии; потенциальному пользователю сессия не создаётся.
\item \textbf{FR-05.} Просмотреть каталог и карточки предметов.
\item \textbf{FR-06.} Создать карточку принесённой вещи с названием, категорией, состоянием и правилами; система должна присвоить экземпляру идентификатор.
\item \textbf{FR-07.} Подтвердить размещение вещи в общей комнате для публикации её карточки.
\item \textbf{FR-08.} Найти предметы по названию.
\item \textbf{FR-09.} Отфильтровать каталог по категории.
\item \textbf{FR-10.} Отфильтровать предметы по доступности на выбранный интервал.
\item \textbf{FR-11.} Автоматически забронировать предмет на интервал без пересечений с его активными бронями; начало не в прошлом, конец позже начала.
\item \textbf{FR-12.} Отменить собственную бронь до получения предмета с освобождением её интервала.
\item \textbf{FR-13.} Отметить получение по своей брони в её интервале, если предмет доступен и предыдущее использование завершено; система должна сохранить время и статус «Используется».
\item \textbf{FR-14.} Отметить возврат с фото предмета в общей комнате: JPEG или PNG до 5 МБ; использование завершается только после сохранения фото.
\item \textbf{FR-15.} Просмотреть собственные брони с фильтром по статусу.
\item \textbf{FR-16.} Просмотреть историю использования предмета с фотографиями возврата при активном подтверждённом доступе.
\item \textbf{FR-17.} Получить уведомления внутри приложения о создании, отмене и просрочке своей брони.
\item \textbf{FR-18.} Отправить сообщение о неисправности или отсутствии предмета; система должна временно запретить его бронирование и получение.
\end{itemize}
\par\noindent\textbf{Администратор должен иметь возможность:}
\begin{itemize}[itemsep=0pt]
\item \textbf{FR-19.} Подтвердить доступ потенциального пользователя: перевести его учётную запись из «Ожидает подтверждения» в «Активен» без повторной регистрации. До подтверждения права пользователя не предоставляются.
\item \textbf{FR-20.} Заблокировать пользователя с отменой неиспользованных броней; возврат уже взятых вещей должен оставаться доступным.
\item \textbf{FR-21.} Отредактировать сведения в карточке предмета.
\item \textbf{FR-22.} Рассмотреть сообщение о проблеме и сохранить решение по нему.
\item \textbf{FR-23.} Изменить доступность предмета; при запрете использования система должна отменить его не начатые брони с причиной.
\end{itemize}
\par\noindent\textbf{Правила выполнения пользовательских операций:}
\begin{itemize}[itemsep=0pt]
\item \textbf{FR-24.} Система должна переводить неиспользованную бронь в статус «Истекло» после окончания её интервала.
\item \textbf{FR-25.} Система должна запрещать новое бронирование и получение предмета, пока его просроченное использование не завершено возвратом.
\item \textbf{FR-26.} Система должна сохранять каждое изменение предмета или брони в истории с временем и исполнителем.
\item \textbf{FR-27.} Система должна хранить пароли как стойкие адаптивные хеши.
\item \textbf{FR-28.} Система должна проверять текущее состояние доступа при каждой защищённой операции, в том числе в уже открытой сессии. После блокировки ранее подтверждённой записи доступны только публичный каталог, вход, выход, просмотр собственных незавершённых использований и их возврат с фото. Блокировка ожидающей подтверждения записи не разрешает вход. Блокировка не переводит человека в потенциального пользователя и не даёт повторно зарегистрировать существующую запись.
\end{itemize}
\clearpage
\setstretch{1.1}
\subsection{U. Удобство использования}
\begin{itemize}[itemsep=3pt]
\item \textbf{U-01.} Интерфейс должен отображать подписи и сообщения на русском языке.
\item \textbf{U-02.} Основные сценарии должны выполняться при ширине экрана от 360 до 1920 пикселей без горизонтальной прокрутки.
\item \textbf{U-03.} При отказе система должна показывать конкретную причину и указывать ошибочное поле, если отказ вызван вводом.
\item \textbf{U-04.} Статусы должны обозначаться текстом, а не только цветом.
\end{itemize}

\subsection{R. Надёжность}
\begin{itemize}[itemsep=3pt]
\item \textbf{R-01.} После перезапуска система должна сохранять записанные предметы, брони, историю и фотографии.
\item \textbf{R-02.} Из 20 одновременных запросов на один свободный интервал система должна успешно выполнить ровно один.
\item \textbf{R-03.} Повтор запроса с тем же идентификатором операции не должен создавать повторную бронь или отметку использования.
\item \textbf{R-04.} Система должна поддерживать восстановление резервной копии PostgreSQL с фотографиями в тестовую среду не более чем за 30 минут.
\item \textbf{R-05.} При ошибке записи в БД система должна откатывать изменения соответствующей операции целиком.
\end{itemize}

\subsection{P. Производительность}
\begin{itemize}[itemsep=3pt]
\item \textbf{P-01.} Система должна обслуживать 50 одновременно работающих пользователей при 1000 предметах и 10\,000 бронях без серверных ошибок, вызванных нагрузкой.
\item \textbf{P-02.} При нагрузке P-01 сервер должен выполнять не менее 95\% запросов каталога и бронирования за 2 секунды, без учёта сетевых задержек и передачи файлов.
\end{itemize}

\subsection{S. Сопровождаемость}
\begin{itemize}[itemsep=3pt]
\item \textbf{S-01.} Приложение должно запускаться на Helios из JAR-файла без Docker и отдельного Node.js-сервера.
\item \textbf{S-02.} Подключение к БД должно настраиваться без изменения исходного кода.
\item \textbf{S-03.} Изменения схемы БД должны применяться версионируемыми миграциями.
\item \textbf{S-04.} Журнал ошибок должен содержать идентификатор запроса и причину сбоя без паролей и секретов.
\end{itemize}
\endgroup
\clearpage
\section{Прецеденты}
Акторы: гость, потенциальный пользователь, зарегистрированный пользователь и администратор. Гость участвует только в регистрации и просмотре каталога. Потенциальный пользователь участвует только в просмотре каталога без авторизации; вход и остальные пользовательские прецеденты ему недоступны. Это отдельный актор без наследования прав гостя или зарегистрированного пользователя. Зарегистрированный пользователь имеет подтверждённый доступ и также может просматривать каталог. Администратор наследует возможности зарегистрированного пользователя. Выход доступен всем вошедшим акторам. Заблокированный пользователь сохраняет только ограниченный доступ к своим незавершённым использованиям и возврату. Загрузка фото обязательна при возврате и показана отношением \texttt{include}.

\begin{figure}[htbp]
\centering
\begin{tikzpicture}[x=0.80cm,y=0.82cm,uc/.style={ellipse,draw,align=center,inner sep=2pt,text width=3.1cm,minimum height=0.85cm,font=\scriptsize\linespread{1}\selectfont},>=Stealth]
\draw (2,0.2) rectangle (14,15.9);
\node[font=\footnotesize\bfseries] at (8,0.75) {Сервис бронирования общих предметов};
\actor{0}{15}{Гость}
\actor{11}{17.5}{Потенциальный\\пользователь}
\actor{0}{7}{Пользователь\\(подтверждённый)}
\actor{16}{10}{Администратор}
\node[uc] (reg) at (5,15) {UC-01.\\Зарегистрироваться};
\node[uc] (catalog) at (5,13.5) {UC-02. Просмотреть\\каталог};
\node[uc] (login) at (11,13.5) {UC-03. Войти\\в систему};
\node[uc] (add) at (5,11.5) {UC-04. Добавить\\предмет};
\node[uc] (book) at (5,10) {UC-05.\\Забронировать};
\node[uc] (cancel) at (5,8.5) {UC-06. Отменить\\бронь};
\node[uc] (take) at (5,7) {UC-07. Отметить\\получение};
\node[uc] (return) at (5,5.5) {UC-08. Отметить\\возврат};
\node[uc] (report) at (5,4) {UC-09. Сообщить\\о проблеме};
\node[uc] (cabinet) at (5,2.5) {UC-10. Просмотреть\\личный кабинет};
\node[uc] (access) at (11,11.5) {UC-11. Управлять\\доступом};
\node[uc] (edit) at (11,10) {UC-12. Изменить\\карточку};
\node[uc] (resolve) at (11,8.5) {UC-13. Рассмотреть\\проблему};
\node[uc] (photo) at (11,5.5) {Загрузить фото\\возврата};
\draw (0.3,14.85)--(reg.west);
\draw (0.3,14.85)--(1,14.85)--(2.5,14.25)--(catalog.north west);
\draw (10.7,17.35)--(8,17.35)--(8,13.5)--(catalog.east);
\foreach \n in {add,book,cancel,take,return,report,cabinet} {\draw (0.3,6.85)--(1.4,6.85)--(\n.west);}
\draw (0.3,6.85)--(1.4,6.85)--(1.4,13.5)--(catalog.west);
\draw (1.4,12.4)--(9,12.4)--(login.south west);
\foreach \n in {access,edit,resolve} {\draw (15.7,9.85)--(14.6,9.85)--(\n.east);}
\draw[dashed,->] (return.east)--node[above=6pt,font=\scriptsize]{\texttt{<<include>>}}(photo.west);
\draw[-{Triangle[open,length=3mm,width=3mm]}] (16,9.35)--(16,-0.8)--(-1.5,-0.8)--(-1.5,6.85)--(-0.3,6.85);
\end{tikzpicture}
\caption{Основные прецеденты и четыре актора системы}
\end{figure}
\clearpage
\newcommand{\uctable}[9]{%
\noindent\begin{minipage}{\textwidth}
\small\setstretch{1.03}\renewcommand{\arraystretch}{1.12}
\begin{tabular}{|P{0.28\textwidth}|P{0.66\textwidth}|}
\hline
Прецедент & UC-#1 \\\hline
ID & #1 \\\hline
Краткое описание & #2 \\\hline
Главный актор & #3 \\\hline
Второстепенные акторы & #4 \\\hline
Предусловия & #5 \\\hline
Основной поток & #6 \\\hline
Альтернативный поток & #7 \\\hline
Постусловия & #8 \\\hline
\end{tabular}
\refstepcounter{table}
\par\smallskip\centering Таблица \thetable{} --- #9.
\end{minipage}\par\medskip
}
\uctable{01}{Гость создаёт учётную запись.}{Гость}{Нет}{У посетителя нет учётной записи.}{1. Гость открывает форму и вводит имя, логин и пароль. 2. Система проверяет данные. 3. Система создаёт учётную запись, ожидающую подтверждения.}{При занятом логине или неверных данных система указывает ошибку; запись не создаётся.}{Создана запись «Ожидает подтверждения»; гость стал потенциальным пользователем. Доступен только каталог, повторная регистрация запрещена. После одобрения новая запись не создаётся.}{Регистрация}
\uctable{02}{Посетитель просматривает доступные предметы.}{Гость; потенциальный пользователь; зарегистрированный пользователь}{Нет}{Сервис доступен в браузере.}{1. Посетитель открывает каталог. 2. Система показывает предметы. 3. Посетитель открывает карточку с описанием, состоянием и занятыми интервалами.}{Пустой каталог отображается без ошибки. Гостю и потенциальному пользователю не показываются история пользователей и фотографии возврата; также не показываются личности владельцев занятых интервалов.}{Посетитель получил информацию о предметах; бронь не создана.}{Просмотр каталога}
\uctable{03}{Пользователь открывает сессию своей учётной записи.}{Зарегистрированный пользователь; администратор}{Нет}{Учётная запись подтверждена администратором; сессия не открыта. Вход разрешён также после блокировки ранее подтверждённой записи, с ограничениями FR-28.}{1. Пользователь вводит логин и пароль. 2. Система проверяет данные и состояние подтверждения записи. 3. Система открывает сессию с разрешениями текущего состояния доступа.}{Неверные данные отклоняются. До подтверждения вход отклоняется с сообщением об ожидании одобрения; сессия не создаётся, доступен только публичный каталог. После блокировки доступны каталог, свои незавершённые использования и их возврат. Права активного пользователя не выдаются.}{При успешном входе создана сессия с разрешениями текущего состояния учётной записи; для ожидающей подтверждения записи сессия отсутствует.}{Вход в систему}
\clearpage
\uctable{04}{Пользователь пополняет общий фонд вещью.}{Зарегистрированный пользователь}{Нет}{Доступ активен; вещь принесена в общую комнату.}{1. Пользователь заполняет карточку. 2. Система присваивает идентификатор. 3. Пользователь маркирует вещь и подтверждает размещение. 4. Система публикует карточку.}{При ошибках полей или без подтверждения размещения карточка не публикуется.}{Предмет включён в общий фонд и доступен в каталоге на общих условиях.}{Добавление предмета}
\uctable{05}{Пользователь резервирует время использования предмета.}{Зарегистрированный пользователь}{Нет}{Доступ активен; предмет доступен, просроченного использования нет.}{1. Пользователь выбирает интервал. 2. Система проверяет время и активные брони. 3. Система автоматически резервирует интервал и показывает бронь в кабинете.}{Неверное время, пересечение или недоступность приводят к отказу с причиной; новая бронь не создаётся.}{Создана бронь «Забронировано»; использование ещё не началось.}{Бронирование}
\uctable{06}{Пользователь освобождает время своей брони.}{Зарегистрированный пользователь}{Нет}{Своя бронь имеет статус «Забронировано».}{1. Пользователь выбирает свою бронь. 2. Нажимает «Отменить». 3. Система проверяет состояние, отменяет бронь и освобождает интервал.}{После получения предмета отмена запрещена; чужая бронь недоступна для изменения.}{Статус --- «Отменено»; запись сохраняется в истории.}{Отмена брони}
\clearpage
\uctable{07}{Пользователь отмечает, что взял вещь из общей комнаты.}{Зарегистрированный пользователь}{Нет}{Доступ активен; наступил интервал своей брони; предмет доступен, предыдущее использование завершено.}{1. Пользователь берёт вещь из общей комнаты. 2. Нажимает «Я взял вещь». 3. Система проверяет условия и сохраняет время получения.}{При отсутствии вещи пользователь сообщает о проблеме по UC-09. При нарушении условий отметка отклоняется.}{Бронь имеет статус «Используется»; у предмета одно незавершённое использование.}{Получение предмета}
\uctable{08}{Пользователь фиксирует возврат вещи с фотографией.}{Зарегистрированный пользователь}{Нет}{Предмет используется по своей брони; возврат разрешён и после блокировки.}{1. Пользователь возвращает вещь в общую комнату. 2. Нажимает «Я вернул вещь» и прикрепляет фото. 3. Система сохраняет фото и время возврата.}{При ошибке формата, размера или загрузки использование не завершается. При повреждении вещи и активном доступе пользователь также выполняет UC-09; после блокировки этот прецедент недоступен, но возврат разрешён.}{Бронь завершена, фотография связана с возвратом. Одобрение другого человека не требуется.}{Возврат предмета}
\uctable{09}{Пользователь сообщает об отсутствии или неисправности вещи.}{Зарегистрированный пользователь}{Администратор --- получатель сообщения}{Доступ активен; выбран предмет.}{1. Пользователь описывает неисправность или отсутствие вещи. 2. Система сохраняет сообщение. 3. Система запрещает использование предмета, отменяет не начатые брони и уведомляет администратора.}{Пустое описание не принимается. Уже начатое использование сохраняется до возврата.}{Сообщение ожидает рассмотрения; предмет недоступен для бронирования и получения.}{Сообщение о проблеме}
\clearpage
\uctable{10}{Пользователь просматривает свои брони и уведомления.}{Зарегистрированный пользователь}{Нет}{Пользователь вошёл; доступ активен либо заблокирован при наличии незавершённого использования.}{1. Пользователь открывает кабинет. 2. Выбирает статус броней. 3. Просматривает свои записи, уведомления и фотографии завершённых возвратов.}{При отсутствии записей показывается пустой список. Чужие брони не доступны для изменения. При блокировке показываются только свои незавершённые использования с действием возврата; история и остальные функции кабинета недоступны. Потенциальному пользователю кабинет недоступен.}{Пользователь получил сведения о своих бронях и возвратах.}{Просмотр личного кабинета}
\uctable{11}{Администратор подтверждает потенциального пользователя или блокирует учётную запись.}{Администратор}{Нет}{Администратор вошёл в систему.}{1. Администратор выбирает учётную запись. 2. Для ожидающей записи подтверждает доступ либо блокирует её с причиной; активную запись может заблокировать с причиной. 3. Система сохраняет новое состояние той же записи; при блокировке отменяет неиспользованные брони.}{Запрос от другой роли отклоняется. Повторное подтверждение активной записи не меняет её и не создаёт дубликат. Блокировка не запрещает возврат уже взятой вещи и не возвращает запись в состояние ожидания.}{Состояние доступа обновлено; действие записано в историю.}{Управление доступом}
\uctable{12}{Администратор исправляет сведения о предмете.}{Администратор}{Нет}{Администратор вошёл; карточка существует.}{1. Администратор открывает карточку. 2. Меняет сведения. 3. Система проверяет поля и сохраняет изменения.}{При ошибке обязательных полей изменения не применяются; система указывает ошибку.}{Сведения карточки обновлены, история изменения сохранена.}{Редактирование карточки}
\clearpage
\uctable{13}{Администратор сохраняет решение по сообщению о проблеме.}{Администратор}{Нет}{Администратор вошёл; сообщение о проблеме существует.}{1. Администратор открывает сообщение. 2. Проверяет ситуацию и записывает решение. 3. После устранения проблемы восстанавливает доступность предмета.}{Если проблема не устранена, предмет остаётся недоступным. Восстановление доступности не возобновляет отменённые брони.}{Решение сохранено, состояние предмета соответствует результату рассмотрения.}{Рассмотрение проблемы}
\clearpage
\section{Архитектура}
Мы планируем придерживаться минималистичного подхода: реализовать сервис как монолитное приложение с понятной структурой и небольшим числом компонентов. Это упростит разработку, развёртывание и сопровождение системы.

Для интерфейса хотим использовать React, для серверной логики --- Spring Boot. Сервер будет обрабатывать регистрацию, проверять права доступа и управлять каталогом, бронированием и возвратами.

В PostgreSQL планируем хранить учётные записи, предметы, брони и историю операций. Дополнительно хотим использовать SeaweedFS как файловое хранилище для фотографий возврата; в PostgreSQL будут храниться ссылки на файлы. Возврат будет завершаться только после успешного сохранения фотографии и записи в БД.
\end{document}
